\chapter{System Design}
\label{ch:4}


\begin{figure}[H]
\centering
\includegraphics[width=0.85\textwidth]{01_system_architecture.png}
\caption{High-Level System Architecture}
\label{fig:01_system_architecture}
\end{figure}


\section{4.1 Overall System Architecture}

\subsection{4.1.1 High-Level System Diagram}
The system comprises three independent modules operating in sequence, with federated learning applied exclusively to Module 1:

See also: 01\_system\_architecture.png in /diagrams/rendered/


\subsection{4.1.2 Federated Learning Communication Flow}
The federated learning component (Module 1) follows a standard synchronous FL protocol:

See also: 04\_fl\_sequence.png (Mermaid Sequence Diagram) in /diagrams/rendered/


\begin{figure}[H]
\centering
\includegraphics[width=0.9\textwidth]{12_class_diagram.png}
\caption{UML Class Diagram}
\label{fig:12_class_diagram}
\end{figure}


\section{4.2 Module 1: Federated Classification Architecture}

\subsection{4.2.1 Client-Server Architecture}
Role: Central Server

\begin{itemize}
  \item Maintains global ResNet-18 model (11.2M parameters)
  \item Coordinates 100 federated rounds
  \item Implements aggregation strategy (FedAvg, FedProx, or QPSO)
  \item Evaluates on balanced global test set
  \item Logs metrics and checkpoints
\end{itemize}

Role: Hospital Clients (3 instances)

\begin{itemize}
  \item Receive broadcast global model from server
  \item Perform local training on private hospital data (5 epochs per round)
  \item Send updated model weights back to server
  \item Receive updated global model for next round
  \item Compute local validation metrics (tracking fairness)
\end{itemize}


\subsection{4.2.2 Model Architecture: ResNet-18 for Brain Tumor Classification}
INPUT (224 × 224 × 3)
  ↓
[Conv 7×7, 64 filters, stride 2]
[BatchNorm + ReLU]
[MaxPool 3×3, stride 2]
  ↓
[ResNet Block 1: 64 filters × 2 blocks]
  ↓
[ResNet Block 2: 128 filters × 2 blocks]
  ↓
[ResNet Block 3: 256 filters × 2 blocks]
  ↓
[ResNet Block 4: 512 filters × 2 blocks]
  ↓
[GlobalAvgPool → 512-dim feature vector]
  ↓
[FC Layer: 512 → 3 classes]
  (Glioma, Meningioma, Pituitary)
  
OUTPUT: Logits (3-dim)
  ↓
[Softmax + Cross-Entropy Loss]
  ↓
Predicted Class + Confidence

Configuration:

\begin{itemize}
  \item **Pretrained Weights:** ImageNet (transfer learning)
  \item **Total Parameters:** 11,178,051 (all trainable)
  \item **Memory:** \~45 MB model size
  \item **Input:** 224×224 RGB images
  \item **Output:** 3-class or 4-class probabilities (depending on experimental setup)
\end{itemize}


\subsection{4.2.3 Aggregation Strategies}

\subsubsection{**Strategy 1: Federated Averaging (FedAvg)**}
\# Pseudocode
for round t in 1..100:
    \# Server broadcasts
    send w\_global to all clients
    
    \# Clients train locally
    for each client c:
        w\_c ← local\_training(data\_c, w\_global, epochs=5)
        send w\_c to server
    
    \# Server aggregates (weighted by dataset size)
    n\_total ← sum of all client dataset sizes
    w\_global ← (0, 0, 0)
    for each client c:
        w\_global ← w\_global + (n\_c / n\_total) * w\_c

Characteristics:

\begin{itemize}
  \item **Deterministic:** Same input always produces same output
  \item **Fast:** O(n\_clients) computation, minimal memory
  \item **Non-memory:** No tracking of previous solutions
  \item **Prone to averaging:** May lose distinct features learned by individual clients
\end{itemize}


\subsubsection{**Strategy 2: Federated Proximal (FedProx)**}
\# Server side: SAME as FedAvg
\# Client side: MODIFIED local training objective

for round t in 1..100:
    send w\_global to all clients
    
    for each client c:
        \# FedProx modification: Add proximal term
        for epoch in 1..5:
            for batch (x, y) in train\_loader\_c:
                pred ← model(x)
                loss\_data ← CrossEntropy(pred, y)
                
                \# Proximal regularization prevents drift
                loss\_prox ← (μ / 2) * ||w\_local - w\_global||²
                loss\_total ← loss\_data + loss\_prox
                
                optimizer.step(loss\_total)
        
        w\_c ← model.state\_dict()
        send w\_c to server
    
    w\_global ← weighted\_average(w\_1, w\_2, w\_3)

Characteristics:

\begin{itemize}
  \item **Regularized:** Proximal term penalizes large client divergence
  \item **μ = 0.01:** Tunable hyperparameter (higher = more constraint)
  \item **Intent:** Stabilize learning under non-IID data
  \item **Risk:** Can over-constrain, preventing beneficial local adaptation
\end{itemize}


\subsubsection{**Strategy 3: Quantum Particle Swarm Optimization FL (QPSO-FL)** - Novel Contribution}
\# Server-side aggregation strategy using QPSO

class QPSOServer:
    def \_\_init\_\_(self, clients, beta=0.7):
        self.personal\_best = \{\}        \# per-client best weights found
        self.personal\_best\_scores = \{\} \# per-client best validation accuracy
        self.global\_best = None        \# best weights across all clients
        self.global\_best\_score = 0.0
        self.mean\_best = None          \# element-wise mean of pbests
        
    def initialize\_particles(self):
        """Initialize all pbests and gbest to current global model"""
        for client in self.clients:
            self.personal\_best[client.id] = deepcopy(self.global\_model)
            self.personal\_best\_scores[client.id] = 0.0
        self.global\_best = deepcopy(self.global\_model)
    
    def aggregate\_qpso(self, client\_weights, client\_val\_accs):
        """
        QPSO Aggregation Step:
        1. Update personal best (pbest) and global best (gbest)
        2. Compute mean best (mbest)
        3. Apply quantum position update
        """
        
        \# Step 1: Update pbest and gbest
        for client\_id, weights in client\_weights.items():
            val\_acc = client\_val\_accs[client\_id]
            
            if val\_acc > self.personal\_best\_scores[client\_id]:
                self.personal\_best[client\_id] = weights
                self.personal\_best\_scores[client\_id] = val\_acc
            
            if val\_acc > self.global\_best\_score:
                self.global\_best = weights
                self.global\_best\_score = val\_acc
        
        \# Step 2: Compute mean best (element-wise mean of all pbests)
        self.mean\_best = compute\_mean\_state\_dict(
            [self.personal\_best[c\_id] for c\_id in self.clients]
        )
        
        \# Step 3: Quantum position update for each parameter
        for param\_name in self.global\_model.state\_dict():
            \# Per-parameter vectors
            φ = random.uniform(0, 1)  \# per-client attraction weight
            u = random.uniform(0.3, 1.0)  \# quantum randomness
            
            \# Attraction point: weighted combination of pbest and gbest
            p = φ * mean\_best[param\_name] + (1 - φ) * gbest[param\_name]
            
            \# Quantum update: stochastic jump with bounded magnitude
            perturbation = β * abs(mean\_best[param\_name] - gbest[param\_name]) * ln(1/u)
            perturbation = clip(perturbation, -0.1, 0.1)  \# Safety bounds
            
            new\_position = p + random.choice([-1, 1]) * perturbation
            
            w\_new[param\_name] = new\_position
        
        return w\_new

QPSO Advantages:

\begin{itemize}
  \item **Memory-Based:** Tracks pbest (personal best) and gbest (global best)
  \item **Stochastic Exploration:** Quantum jumps escape local optima
  \item **Adaptive:** Attraction point balances individual and global solutions
  \item **Fairness:** Validates based on loss surface (implicitly penalizes unfair models)
\end{itemize}

Parameters:

\begin{itemize}
  \item **β = 0.7:** Contraction-expansion coefficient (exploration vs. exploitation)
  \item **u ∈ [0.3, 1.0]:** Quantum randomness range
  \item **Perturbation clamp [-0.1, 0.1]:** Safety constraint (prevents divergence)
\end{itemize}

See also: 05\_aggregation\_strategies.png (comparison of all 3 strategies)


\begin{figure}[H]
\centering
\includegraphics[width=0.85\textwidth]{13_use_case_diagram.png}
\caption{UML Use Case Diagram}
\label{fig:13_use_case_diagram}
\end{figure}


\section{4.3 Module 2: 3D Brain Tumor Segmentation Architecture}

\subsection{4.3.1 Model Architecture: 3D Attention U-Net}
Attention Gate Mechanism (per skip connection):

Configuration:

\begin{itemize}
  \item **Channels:** 16, 32, 64, 128, 256 (per MONAI standard)
  \item **Strides:** (2, 2, 2, 2) depth-wise
  \item **Patch Size (during training):** 96³ voxels
  \item **Loss Function:** Generalized Dice Loss (handles class imbalance)
  \item **Optimizer:** Adam (lr=1e-4, weight\_decay=1e-5)
  \item **Training Epochs:** 20
  \item **Batch Size:** 1 (volumetric processing)
\end{itemize}

See also: 03\_unet\_architecture.png (detailed layer structure)


\subsection{4.3.2 Preprocessing Pipeline (MONAI)}
Raw NIfTI File (multi-modality)
  ↓
[LoadImage] → Load all 4 modalities separately
  ↓
[Orientation] → Reorient to RAS (Right-Anterior-Superior) standard
  ↓
[Spacing] → Resample to 1.0 × 1.0 × 1.0 mm³ isotropic
  ↓
[NormalizeIntensity] → Z-score normalization (nonzero voxels only)
  ↓
[RandCropByPosNegLabel] → Extract 96³ patches (balanced pos/neg samples)
  ↓
Preprocessed Tensor (1, 4, 96, 96, 96)
  ↓
[Model Forward Pass]
  ↓
Segmentation Output (1, 3, 96, 96, 96)
  ↓
[Inverse Transforms] → Map back to original space
  ↓
Final Segmentation Mask (same space as input)


\begin{figure}[H]
\centering
\includegraphics[width=0.9\textwidth]{14_activity_diagram.png}
\caption{UML Activity Diagram}
\label{fig:14_activity_diagram}
\end{figure}


\section{4.4 Module 3: Progression Forecasting Architecture}

\subsection{4.4.1 Mathematical Models Pipeline}
Grade-Stratified Fitting:

\begin{itemize}
  \item **HGG (High-Grade):** 89 patients, aggressive growth (exponential-like)
  \item **LGG (Low-Grade):** 22 patients, slow growth (logistic/gompertz)
\end{itemize}


\subsection{4.4.2 LSTM Hybrid Architecture}
Loss Function: Mean Squared Error (MSE) on residuals

Training Configuration:

\begin{itemize}
  \item **Optimizer:** Adam (lr=0.001, weight\_decay=1e-5)
  \item **Learning Rate Scheduler:** ReduceLROnPlateau (factor=0.5, patience=10)
  \item **Epochs:** 50-100 (until convergence)
  \item **Batch Size:** 8 (sequences per batch)
  \item **Train/Test Split:** 80/20 per grade
\end{itemize}


\subsection{4.4.3 Prediction Pipeline}
See also: 07\_progression\_pipeline.png (detailed flowchart)


\begin{figure}[H]
\centering
\includegraphics[width=0.85\textwidth]{15_component_diagram.png}
\caption{UML Component Diagram}
\label{fig:15_component_diagram}
\end{figure}


\section{4.5 Cross-Module Data Flow}

\subsection{4.5.1 Complete End-to-End Data Flow}

\begin{figure}[H]
\centering
\includegraphics[width=0.85\textwidth]{16_deployment_diagram.png}
\caption{UML Deployment Diagram}
\label{fig:16_deployment_diagram}
\end{figure}


\section{4.6 Error Handling and Data Validation}

\subsection{4.6.1 Input Validation (All Modules)}
Classification Module:

\begin{itemize}
  \item Check image dimensions: must be 224×224
  \item Check color channels: must be 3 (RGB)
  \item Check pixel range: must be [0, 1] after normalization
  \item Reject images with >10\% black borders (non-brain)
\end{itemize}

Segmentation Module:

\begin{itemize}
  \item Check 3D volume dimensions: typically 155×188×155 (BraTS standard)
  \item Check 4 modalities present: T1, T1ce, T2, FLAIR
  \item Check pixel range after preprocessing: typically [-3, 3] (Z-norm)
  \item Reject volumes with >50\% missing slices
\end{itemize}

Progression Module:

\begin{itemize}
  \item Check ≥ 2 timepoints per patient (minimum for fitting)
  \item Check temporal ordering: timepoints must be chronologically ordered
  \item Check volume continuity: no abrupt 10x jumps (likely segmentation error)
  \item Flag suspicious outliers for manual review
\end{itemize}


\subsection{4.6.2 Model Robustness}
Federated Learning:

\begin{itemize}
  \item If 1 client unavailable: Continue with 2 clients (degraded fairness evaluation)
  \item If validation loss diverges: Early stopping at round t
  \item Model checkpointing: Save best model per 10 rounds
\end{itemize}

Segmentation:

\begin{itemize}
  \item Catch GPU out-of-memory: Reduce patch size dynamically
  \item Handle 1-2 modality inputs: Pad missing modalities with zeros (with warning)
\end{itemize}

Progression:

\begin{itemize}
  \item Handle missing timepoints: Interpolate (with uncertainty flag)
  \item Handle contradictory models: Use ensemble prediction (average all model predictions)
\end{itemize}


\begin{figure}[H]
\centering
\includegraphics[width=0.9\textwidth]{17_state_chart.png}
\caption{UML State Chart Diagram --- FL Training Lifecycle}
\label{fig:17_state_chart}
\end{figure}


\section{4.7 System Configuration and Hyperparameters}

\subsection{4.7.1 Module 1: Federated Classification}

\subsection{4.7.2 Module 2: Segmentation}

\subsection{4.7.3 Module 3: Progression}

\begin{figure}[H]
\centering
\includegraphics[width=0.8\textwidth]{04_fl_sequence.png}
\caption{FL Communication Sequence Diagram}
\label{fig:04_fl_sequence}
\end{figure}


\section{4.8 Chapter Summary}
Figure: Federated Learning System Architecture

Figure: 3D U-Net Segmentation Architecture

Figure: Aggregation Strategies Comparison

Figure: End-to-End Integration Architecture


\section{Chapter Summary}
This chapter detailed the system design with UML diagrams, data flow diagrams, and mathematical formulations for all three aggregation strategies.
